\chapter{化学符号急救箱}

\begin{kaichang}
想象你正在读这本书的某一章，读到一句话：“\chem{2H_2O_2} 在过氧化氢酶的催化下分解成 \chem{2H_2O} 和 \chem{O_2}。”你的脑子突然卡住了：那个右下角的小 2 是什么意思？为什么前面又站着一个大 2？箭头右边那堆字母组合凭什么等于左边那堆？

别慌，这不是你笨，是符号系统在跟你打暗号。化学符号是一套压缩率极高的“速写语言”：一个式子里藏着成分、数量、连接方式和变化方向四层信息。正因为它压缩得厉害，第一次解压失败再正常不过。

这份附录就是你的解压工具箱。它不按故事顺序讲，而是按“你卡在哪里”来组织：卡在元素符号就查第一小节，卡在化学式就查第二小节，卡在配平、摩尔、浓度、pH 的计算就直接翻到最后的分步方法清单。建议你先通读一遍混个脸熟，然后把它留在手边——读正文卡壳时，随时回来急救。
\end{kaichang}

\section{这套符号是怎么来的}

在动手查表之前，先花一分钟知道这套系统的“设计思路”，后面所有规则都会显得顺理成章。

化学符号要解决一个难题：全世界说不同语言的人，怎样谈论同一堆看不见的原子而不吵架？答案是用一套拉丁字母做“世界语”。每种元素分到一两个字母，称为元素符号；把符号按一定规则拼起来，就得到化学式；在化学式之间画上箭头，就得到化学方程式。从符号到方程式，信息一层比一层多，就像从字母到单词再到句子。

还有一条全书通用的约定要记牢：符号只负责“记账”，不负责“画画”。\chem{H_2O} 里的字母和数字只告诉你“两个氢原子加一个氧原子结成一伙”，它\textbf{不}告诉你这伙人长什么样、弯成多少度、彼此挨多近。想知道形状，要靠结构式和模型（附录 C 专门讲各种分子模型各自保留了什么、牺牲了什么）。把符号当成账本而不是照片，你就不会被它误导。

你可能想问：这套速写是谁定的？元素符号的通行方案出自瑞典化学家贝采利乌斯（1813 年前后），他的理由很务实——当时的化学家各画各的图形符号，一篇文章换一个作者就得重学一套，而用拉丁字母做代号，全世界认字的人都能读写。两百年来这套系统只打过补丁、没换过地基，你今天学的就是全世界实验室通用的那一套。学会它，你就拿到了和任何国家的化学家“笔谈”的资格。

\section{元素符号：化学的字母}

元素符号的规矩只有三条：

\begin{enumerate}[itemsep=2pt]
\item 一个字母的，大写：\chem{H}（氢）、\chem{O}（氧）、\chem{C}（碳）。
\item 两个字母的，\textbf{一大一小}：\chem{Ca} 是钙，\chem{CA} 什么也不是，\chem{ca} 也什么也不是。大小写在这里不是装饰，写错就换了一个词。
\item 有些符号来自拉丁名而不是中文或英文名，所以看起来“不讲道理”：铁 \chem{Fe}（ferrum）、钠 \chem{Na}（natrium）、钾 \chem{K}（kalium）、铜 \chem{Cu}（cuprum）、银 \chem{Ag}（argentum）、金 \chem{Au}（aurum）、铅 \chem{Pb}（plumbum）、汞 \chem{Hg}（hydrargyrum）。这几个是历史留下的“钉子户”，只能眼熟，没法推导。
\end{enumerate}

下面这张表收了全书出场最频繁的元素。不需要背，只需要“见过”。

\begin{table}[htbp]
\centering
\caption{全书高频元素速查}
\begin{tabular}{llll}
\toprule
符号 & 名称 & 一句话印象 & 生活里在哪见过 \\
\midrule
\chem{H} & 氢 & 最轻的元素 & 水、酸、恒星燃料 \\
\chem{C} & 碳 & 搭骨架的能手 & 铅笔芯、糖、你 \\
\chem{N} & 氮 & 空气的主角 & 空气约八成是它 \\
\chem{O} & 氧 & 抢电子的活跃分子 & 呼吸、铁锈、水 \\
\chem{P} & 磷 & 生命账本的关键 & 骨骼、ATP、DNA \\
\chem{S} & 硫 & 有特殊气味 & 蛋白质、火山口 \\
\chem{Na} & 钠 & 遇水暴躁的金属 & 食盐（以离子形式） \\
\chem{K} & 钾 & 神经信号的主角 & 香蕉、神经冲动 \\
\chem{Ca} & 钙 & 建筑与信使 & 骨骼、石灰、蛋壳 \\
\chem{Mg} & 镁 & 叶绿素的心脏 & 绿叶、补镁片 \\
\chem{Fe} & 铁 & 变价高手 & 血液、钉子、铁锈 \\
\chem{Cu} & 铜 & 导电的老朋友 & 电线、铜锣烧的锅 \\
\chem{Zn} & 锌 & 酶的常客 & 镀锌铁皮、免疫 \\
\chem{Cl} & 氯 & 抢到电子才安稳 & 食盐、消毒水 \\
\chem{Si} & 硅 & 沙子与芯片 & 玻璃、手机芯片 \\
\chem{I} & 碘 & 甲状腺的原料 & 加碘盐、碘酒 \\
\bottomrule
\end{tabular}
\end{table}

\textbf{急救方法：读一个陌生的元素符号}

\begin{enumerate}[itemsep=2pt]
\item 先看首字母，一定大写；若有两个字母，第二个必是小写。
\item 翻周期表（附录 B）找到它，确认名称和它在表中的位置。
\item 看它所在的列（族）：同一列脾气相近，这能帮你猜它在反应里会干什么。
\item 问自己：此刻它是原子、单质分子还是离子？同一个符号在不同上下文里角色不同——\chem{Cl} 可以指氯元素、一个氯原子，而食盐里实际存在的是氯离子 \chem{Cl^-}。
\end{enumerate}

\section{化学式：一份精确的配料单}

把元素符号拼在一起，加上几个数字，就成了化学式。化学式回答贯穿全书的第一个问题：“它由什么组成，各有多少？”

规则其实只有四条：

\begin{enumerate}[itemsep=2pt]
\item \textbf{下标管左边那一个符号}。\chem{H_2O} 表示“2 个氢原子加 1 个氧原子”：下标 2 只属于 \chem{H}，\chem{O} 没有下标就是 1。一个水分子总共 3 个原子。
\item \textbf{括号管里面一整包}。\chem{Ca(OH)_2} 表示 1 个钙、2 个氧、2 个氢——括号外的下标 2 把括号里的 \chem{OH} 整组复制一份。数原子时先拆括号。
\item \textbf{右上角的正负号是电荷}。\chem{SO_4^{2-}} 是硫酸根离子：1 个硫、4 个氧，整团多带了 2 个电子，所以带 2 个单位负电荷。下标数原子，上标数电荷，两码事，位置也不同。
\item \textbf{式子最前面的大数字是“几份”}。\chem{2H_2O} 是 2 个水分子，共 4 个氢原子、2 个氧原子。这个系数属于整个式子，像购物清单上的“×2”。
\end{enumerate}

\begin{wuqu}
“\chem{H_2O} 里的 2 说明氧有 2 个。”——错。下标永远只管它左边紧挨着的那一个符号。2 属于氢。同样，\chem{CO_2} 是“1 个碳加 2 个氧”，不是“2 个碳加 2 个氧”。下次看到一个化学式，先把每个下标“认领”到它左边的符号，再开始数原子，就不会数错。
\end{wuqu}

还有一类式子容易唬人：带结晶水的写法，比如胆矾 \chem{CuSO_4\cdot 5H_2O}。中间那个点不是乘号，也不表示化学反应，它表示“每 1 份硫酸铜的晶体里，规规矩矩地嵌着 5 个水分子”。加热能把这些水赶走，蓝色晶体变白——水只是住在晶体里的房客，不是硫酸铜的一部分。

\textbf{急救方法：数清一个化学式里的原子}

\begin{enumerate}[itemsep=2pt]
\item 去掉最前面的系数，先只看一份。
\item 从左到右，把每个符号的下标记下来（没写下标的记 1）。
\item 遇到括号，把括号内各原子的下标乘以括号外的下标。
\item 遇到“\chem{\cdot xH_2O}”这类写法，把结晶水单独记一笔。
\item 最后把同种元素的原子数加起来；若有前系数，再整体乘以系数。
\end{enumerate}

练一个：\chem{3Ca(OH)_2}。去掉系数看 \chem{Ca(OH)_2}：钙 1，括号内氧 1、氢 1，括号外下标 2，所以氧 2、氢 2；再乘前系数 3：钙 3、氧 6、氢 6。一共 15 个原子。

顺带分清一个全书反复使用的区别：\textbf{化学式对分子和对离子化合物，含义不完全一样}。对水、二氧化碳这样的分子，化学式忠实记录一个分子里每种原子的真实个数，\chem{H_2O} 就是“一个分子含 2 氢 1 氧”。但对食盐这样的离子化合物，晶体里根本没有“一个氯化钠分子”，只有按 1:1 交替排列、望不到边的钠离子和氯离子，\chem{NaCl} 只表示\textbf{两种离子的最简个数比}。同样，金刚石写作 \chem{C}，石墨也写作 \chem{C}，因为它们是原子连成的大网，没有单个分子可数，只能写成分。看到化学式先问一句“这是分子式还是最简比”，很多困惑会当场消失。

\section{结构式：谁连着谁}

化学式只说“有什么、有多少”，不说“怎么连”。可是第 20 章讲过一个关键事实：分子的形状和连接方式决定它会做什么。葡萄糖和果糖化学式都是 \chem{C_6H_{12}O_6}，甜度却不同——差别全在连接方式上。结构式就是专门表达“谁连着谁”的符号。

最基本的约定：\textbf{一条短线代表一对共享的电子，也就是一个共价键}。水写成 \chem{H-O-H}，氧伸两只手各拉一个氢；二氧化碳写成 \chem{O=C=O}，两条短线表示双键（共享两对电子）；氮气写成 $\chem{N}\equiv\chem{N}$，三条短线是牢固的三键。画结构式时，每种原子能伸出的“手”数基本固定：氢 1 只、氧 2 只、氮 3 只、碳 4 只。这个“手数”就是第 17 到 19 章讲的成键规律在符号上的投影，遇到陌生结构时先数手数，画错了一查就知道。

有机分子太大，全画出来累死人，于是化学家发明了逐级偷懒的写法：

\begin{itemize}[itemsep=2pt]
\item \textbf{完整结构式}：每条键都画，如乙醇画成 \chem{H_3C-CH_2-OH} 的展开版，清楚但占地方。
\item \textbf{简写式}：把连在同一碳上的氢合并，写成 \chem{CH_3CH_2OH}。乙醇、醋酸 \chem{CH_3COOH} 常用这种。
\item \textbf{骨架式}：碳原子不写字母，用折线的端点和拐点代替，氢能省则省。正文里的复杂有机分子和生物分子多用这种。骨架式的完整读法（楔线、立体信息）在附录 C，这里只需记住：\textbf{折线上每个端点和拐点都是一个碳，碳的四只手没画出来的部分默认被氢填满}。
\end{itemize}

为什么要费这个事？看个对比例子就懂。乙醇和二甲醚的化学式都是 \chem{C_2H_6O}——单看配料单，两者一模一样。但写出连接方式：乙醇是 \chem{CH_3-CH_2-OH}，氧插在碳链末端、带着一个氢；二甲醚是 \chem{CH_3-O-CH_3}，氧被夹在两个碳中间。就这一点差别，前者是酒类饮料里能让人微醺的液体（能和水任意混合），后者是常温下容易挥发的气体，两者脾气天差地别。化学式相同的物质未必是同一种物质，\textbf{连接方式也是身份的一部分}——这正是结构式存在的全部理由，也是第 20 章“形状决定功能”和第 42 章手性故事的符号前奏。

\section{反应箭头：化学句子的语法}

化学方程式是用化学式写的句子，而箭头是这句子的谓语。读句子先找谓语，读方程式先看箭头——以及箭头周围的标注。

\begin{table}[htbp]
\centering
\caption{箭头和它的搭档们}
\begin{tabular}{ll}
\toprule
写法 & 意思 \\
\midrule
$\rightarrow$ & 反应朝这个方向进行，左边是反应物，右边是生成物 \\
$\rightleftharpoons$ & 两个方向都在发生，处于动态平衡（第 31 章） \\
箭头上方写字 & 反应条件，如“加热”“光照”“催化剂”或某种酶 \\
$\uparrow$ 标在生成物后 & 该物质以气体形式离开体系 \\
$\downarrow$ 标在生成物后 & 该物质以沉淀形式离开体系 \\
\bottomrule
\end{tabular}
\end{table}

举一句完整的：“\chem{2H_2O_2} $\xrightarrow{\text{过氧化氢酶}}$ \chem{2H_2O} + \chem{O_2}\,$\uparrow$”。读法：两份过氧化氢，在过氧化氢酶的帮助下，变成两份水和一份氧气，氧气冒泡跑掉。加号读“和”，箭头读“变成”，系数读“几份”——一个方程式就这样还原成了人话。

\begin{wuqu}
“箭头两边物质不一样，说明原子被消灭或新生了。”——不对。化学反应里原子不生不灭，只是\textbf{重新排队}（第 26 章把方程式叫做“反应的账本”，就是这个道理）。数一下上面的例子：左边 4 个氢、4 个氧；右边 \chem{2H_2O} 含 4 个氢、2 个氧，\chem{O_2} 含 2 个氧，合计也是 4 氢 4 氧。账平了。任何写出的方程式如果两边原子数对不上，那不是自然界的奇迹，是写字的人还没配平。
\end{wuqu}

还要注意箭头的“语气”。单向箭头不等于“永远只能朝这边走”，它只表示在这套条件下主要朝这边进行；可逆符号 $\rightleftharpoons$ 也不是“停住不动”，而是两个方向同时跑、速率恰好相等（第 31 章）。至于催化剂和酶写在箭头上方，意思是“帮忙搭了条好走的路”，它们\textbf{不出现在账本里}——不被消耗，也不改变最终能变多少，只改变到达的速度（第 30、49 章）。

\textbf{急救方法：读懂一个陌生方程式}

\begin{enumerate}[itemsep=2pt]
\item 先找箭头，分清反应物（左）和生成物（右）。
\item 把每个化学式按上一节的方法拆成原子清单，确认两边每种原子数目相等。
\item 看箭头上下有没有条件：加热、催化剂、酶？这告诉你“快不快、靠什么快”，但改变不了“能不能”。
\item 看有没有 $\uparrow$、$\downarrow$，它们解释了现象：冒泡、变浑浊。
\item 最后问自己一句：这个反应里，旧键断裂、新键生成，能量总体是放出还是收进？方程式本身不写能量，答案要回第 27、28 章找。
\end{enumerate}

\section{常见离子速查表}

离子是“电子账本不平”的原子或原子团：丢了电子带正电（阳离子），多拿了电子带负电（阴离子）。写离子符号时，电荷数标在右上角，数字在前、正负号在后，如 \chem{Ca^{2+}}；电荷为 1 时省略数字，写 \chem{Na^+} 不写 \chem{Na^{1+}}。

\begin{table}[htbp]
\centering
\caption{全书常见阳离子}
\begin{tabular}{lll}
\toprule
符号 & 名称 & 在哪里遇到 \\
\midrule
\chem{H^+} & 氢离子（质子） & 酸的共同主角（第 32 章） \\
\chem{Na^+} & 钠离子 & 食盐、体液、神经信号 \\
\chem{K^+} & 钾离子 & 神经冲动的另一半 \\
\chem{NH_4^+} & 铵根离子 & 化肥、代谢产物 \\
\chem{Ca^{2+}} & 钙离子 & 骨骼、肌肉收缩的信号 \\
\chem{Mg^{2+}} & 镁离子 & 叶绿素中心、酶的助手 \\
\chem{Fe^{2+}} & 亚铁离子 & 血红蛋白抓住氧的那一端 \\
\chem{Fe^{3+}} & 铁离子 & 铁锈里的主角 \\
\chem{Cu^{2+}} & 铜离子 & 蓝色溶液、胆矾 \\
\chem{Zn^{2+}} & 锌离子 & 许多酶的工作位点 \\
\bottomrule
\end{tabular}
\end{table}

\begin{table}[htbp]
\centering
\caption{全书常见阴离子}
\begin{tabular}{lll}
\toprule
符号 & 名称 & 在哪里遇到 \\
\midrule
\chem{Cl^-} & 氯离子 & 食盐、胃酸、体液 \\
\chem{OH^-} & 氢氧根离子 & 碱的共同主角（第 32 章） \\
\chem{O^{2-}} & 氧离子 & 金属氧化物晶体中 \\
\chem{CO_3^{2-}} & 碳酸根离子 & 石灰石、贝壳、苏打 \\
\chem{HCO_3^-} & 碳酸氢根离子 & 小苏打、血液的缓冲液 \\
\chem{SO_4^{2-}} & 硫酸根离子 & 硫酸盐、胆矾 \\
\chem{NO_3^-} & 硝酸根离子 & 化肥、腌菜的争议角色 \\
\chem{PO_4^{3-}} & 磷酸根离子 & 骨骼、ATP、DNA 骨架 \\
\chem{CH_3COO^-} & 醋酸根离子 & 醋的酸味来源之一 \\
\bottomrule
\end{tabular}
\end{table}

三个提醒。第一，\chem{Fe^{2+}} 和 \chem{Fe^{3+}} 是同一种元素铁，只差一个电子，性质却差很多——离子符号上的电荷不是小装饰，是身份的一部分。第二，像 \chem{SO_4^{2-}} 这样好几个原子抱团带一个总电荷的，叫多原子离子（原子团），它们在多数反应里\textbf{整团行动、不拆伙}，配平时可以当成一个整体来数。第三，离子从不单独存在：一瓶盐水里没有纯粹的 \chem{Na^+}，它身边永远围着等量的负电荷，整个体系电中性。写离子方程式或算浓度时，忘了这个约束就会算出“一瓶带电的水”这种不存在的东西。

离子电荷其实常常可以“猜”出来，猜的依据就在周期表的位置（第 10、11 章）。第 1 族的金属最外层只有 1 个电子，丢掉后内层刚好住满，于是总形成 $+1$ 离子；第 2 族类似，形成 $+2$ 离子。第 17 族最外层差 1 个电子住满，抢到 1 个就安稳，于是总形成 $-1$ 离子；第 16 族（氧、硫）差 2 个，常形成 $-2$ 离子。这不是原子有什么愿望，而是这些排布对应的总能量更低（第 11、18 章）。中间街区的过渡金属没这么省事：\chem{Fe^{2+}} 与 \chem{Fe^{3+}}、\chem{Cu^+} 与 \chem{Cu^{2+}} 都稳定存在，它们的变价正是第 15 章和第 33 章的故事素材。记离子不用死背整张表，先记“族号定主价”，再把过渡金属的多变当作特例，负担立刻小了一半。

\section{官能团与生物分子图例}

有机分子动不动几十个原子，但真正决定它“会干什么”的，往往只是身上挂着的几个小零件——官能团（第 37 章把它们叫作“反应接口”）。认出零件，就猜得到大半脾气。

\begin{table}[htbp]
\centering
\caption{常见官能团速查}
\begin{tabular}{llll}
\toprule
写法 & 名称 & 脾气 & 出现在哪 \\
\midrule
\chem{-OH} & 羟基 & 亲水、可形成氢键 & 酒精、糖 \\
\chem{-CHO} & 醛基 & 容易被氧化 & 葡萄糖、香料 \\
\chem{-CO-} & 羰基 & 极性、反应位点 & 丙酮、糖 \\
\chem{-COOH} & 羧基 & 能放出 \chem{H^+}，显酸性 & 醋、柠檬酸 \\
\chem{-NH_2} & 氨基 & 能接住 \chem{H^+}，显碱性 & 氨基酸、蛋白质 \\
\chem{-COO-} & 酯键（酯基） & 香味来源、油脂的缝线 & 果香、脂肪 \\
\chem{-PO_4^{2-}} & 磷酸基 & 带负电、参与能量转账 & ATP、DNA 骨架 \\
\chem{-CH_3} & 甲基 & 惰性标记，可“贴标签” & 表观遗传的记号 \\
\bottomrule
\end{tabular}
\end{table}

正文后半册还会出现一批生物分子的“速写记号”，在这里集中认一遍：

\begin{itemize}[itemsep=2pt]
\item \textbf{葡萄糖}：\chem{C_6H_{12}O_6}。这个化学式值得眼熟——光合作用的产物、呼吸作用的燃料，第六篇的主线人物。
\item \textbf{氨基酸通式}：一个中心碳，连着 \chem{-NH_2}、\chem{-COOH}、\chem{-H} 和一个侧链（常记作 \chem{R}）。二十种氨基酸的差别全在 \chem{R} 上（第 48 章）。
\item \textbf{ATP 与 ADP}：腺苷三磷酸和腺苷二磷酸。写 \chem{ATP} $\rightarrow$ \chem{ADP} + \chem{P_i} 表示 ATP 把末端磷酸基转交给别人。提醒一句：能量并不“藏”在某一根特殊的高能键里，放能是因为产物整体比反应物更稳定（第 50 章专门拆这个比喻）。
\item \textbf{DNA 的四个字母}：A、T、G、C 是四种碱基（腺嘌呤、胸腺嘧啶、鸟嘌呤、胞嘧啶）的缩写；RNA 里 T 换成 U（尿嘧啶）。配对规矩：A 对 T（或 U），G 对 C（第 66 章）。
\item \textbf{氨基酸的单字母码}：测序和基因章节里会见到 A、G、S 等单字母代表氨基酸。注意它们与碱基字母是\textbf{两套不同的字典}，同一个字母 A 在 DNA 序列里是一种碱基、在蛋白质序列里是丙氨酸，看上下文分辨。
\end{itemize}

\section{方法一：配平化学方程式}

配平就是给方程式配上系数，让两边每种原子的数目相等——把账本做平。规则只有一条铁律：\textbf{只许改式子前面的系数，不许改化学式里的下标}。改下标等于换一种物质，账就做成假账了。

以甲烷燃烧为例，骨架是 \chem{CH_4} + \chem{O_2} $\rightarrow$ \chem{CO_2} + \chem{H_2O}：

\begin{enumerate}[itemsep=2pt]
\item \textbf{列清单}。左边：碳 1、氢 4、氧 2。右边：碳 1、氢 2、氧 3。碳已经平了，先不动。
\item \textbf{从只出现一次的元素下手}。氢左边 4、右边 2，给 \chem{H_2O} 配系数 2：右边变成氢 4、氧 $2+2=4$。现在氢平了。
\item \textbf{最后收拾单质}。氧右边共 4，左边 \chem{O_2} 配系数 2 即 4。平了。
\item \textbf{总复核}。\chem{CH_4} + \chem{2O_2} $\rightarrow$ \chem{CO_2} + \chem{2H_2O}：两边都是碳 1、氢 4、氧 4。收工。
\end{enumerate}

偶尔会出现分数：比如配平 \chem{C_2H_6} + \chem{O_2} $\rightarrow$ \chem{CO_2} + \chem{H_2O} 时，氧可能需要 $\frac{7}{2}$。别慌，先留着分数把账配平，最后一步把所有系数同乘 2 消掉分母：\chem{2C_2H_6} + \chem{7O_2} $\rightarrow$ \chem{4CO_2} + \chem{6H_2O}。系数必须是整数，因为分子只能整个整个地反应。

再提醒一次顺序口诀：\textbf{先配原子团复杂的，再配只出现一次的，单质留到最后}。氧、氢常常出现在好几个式子里，先动它们容易顾此失彼；单质（如 \chem{O_2}）只出现在一处，最后用它兜底最省力。

\section{方法二：摩尔换算}

摩尔是化学的“打”——就像鸡蛋按打算，原子和分子按摩尔算，只不过一“摩尔”是 \ud{6.02\times 10^{23}}{} 个（阿伏伽德罗常数，第 5 章讲它怎么测出来）。换算的桥是\textbf{摩尔质量}：数值上等于化学式量，单位是 \ud{}{g/mol}。水的摩尔质量：氢约 \ud{1}{g/mol}，氧约 \ud{16}{g/mol}，所以 \chem{H_2O} 是 $2\times1+16=18\ \ud{}{g/mol}$。

\textbf{三步换算（质量 $\rightarrow$ 摩尔 $\rightarrow$ 粒子数）}

\begin{enumerate}[itemsep=2pt]
\item 写出化学式，算出摩尔质量 $M$（把式中各原子的相对原子质量按下标加总）。
\item 物质的量 $n = \dfrac{m}{M}$：质量除以摩尔质量，单位是 mol。
\item 粒子数 $N = n\times 6.02\times 10^{23}$。反过来算，就倒着走：粒子数除以阿伏伽德罗常数得摩尔，再乘摩尔质量得质量。
\end{enumerate}

例题：\ud{36}{g} 水是多少摩尔、多少个分子？

\begin{enumerate}[itemsep=2pt]
\item $M(\chem{H_2O}) = 18\ \ud{}{g/mol}$。
\item $n = 36 \div 18 = 2\ \ud{}{mol}$。
\item $N = 2\times 6.02\times 10^{23} = 1.2\times 10^{24}$ 个水分子。
\end{enumerate}

一杯 \ud{36}{g} 的水里大约有 $10^{24}$ 个分子——比全地球的沙滩上的沙子总数还多得多。这就是摩尔存在的理由：它把“数不过来的小”和“称得出来的大”接在了一起。

\begin{qiaoliang}
方程式配平后的系数，正是\textbf{摩尔数之比}。\chem{CH_4} + \chem{2O_2} $\rightarrow$ \chem{CO_2} + \chem{2H_2O} 读作：1 mol 甲烷恰好和 2 mol 氧气反应，生成 1 mol 二氧化碳和 2 mol 水。于是 \ud{16}{g}（1 mol）甲烷需要 $2\times32=64$ \ud{}{g} 氧气——工业上算原料用量、实验室里算产物产量，用的都是这一句话。系数不是装饰，是化学的汇率。
\end{qiaoliang}

\section{方法三：浓度计算}

浓度回答“这份溶液里溶了多少东西”。全书用到两种常见说法：

\begin{itemize}[itemsep=2pt]
\item \textbf{质量浓度}：单位体积里有多少克，如 \ud{}{g/L}。生理盐水的“\ud{0.9}{\%}”指每 \ud{100}{mL} 溶液含 \ud{0.9}{g} 氯化钠，即 \ud{9}{g/L}。
\item \textbf{物质的量浓度}：单位体积里有多少摩尔，记作 $c$，单位 \ud{}{mol/L}。公式：$c = \dfrac{n}{V}$，$n$ 是溶质的物质的量，$V$ 是溶液总体积（以升计）。
\end{itemize}

\textbf{三步算浓度}

\begin{enumerate}[itemsep=2pt]
\item 用第二节的方法把溶质质量换算成摩尔：$n = m/M$。
\item 把体积统一换算成升（\ud{250}{mL} = \ud{0.25}{L}）。
\item 代入 $c = n/V$。
\end{enumerate}

例题：把 \ud{5.85}{g} 食盐溶于水配成 \ud{500}{mL} 溶液，浓度是多少？$\chem{NaCl}$ 的摩尔质量约 $23+35.5=58.5\ \ud{}{g/mol}$，所以 $n=5.85\div58.5=0.1\ \ud{}{mol}$；$V=0.5\ \ud{}{L}$；$c=0.1\div0.5=0.2\ \ud{}{mol/L}$。

稀释问题用一句口诀：\textbf{加水只变体积，不变溶质的摩尔数}。所以稀释前后 $c_1V_1 = c_2V_2$。把上面 \ud{0.2}{mol/L} 的盐水取 \ud{100}{mL} 加水到 \ud{400}{mL}，新浓度 $c_2 = 0.2\times100/400 = 0.05\ \ud{}{mol/L}$。医学里配药、农业里配营养液，靠的都是这一行算式——剂量差十倍，后果可能完全不同，所以浓度从来不是小事。

\section{方法四：pH 计算}

pH 是水溶液里氢离子浓度的“对数刻度”（第 32 章讲它为什么这样设计）：

\[
\mathrm{pH} = -\lg\,c(\chem{H^+})
\]

其中 $c(\chem{H^+})$ 以 \ud{}{mol/L} 计。对数刻度的好处是把跨度极大的数字压扁：\chem{H^+} 浓度从 \ud{1}{mol/L} 到 \ud{10^{-14}}{mol/L}，pH 恰好从 0 走到 14。

\textbf{三步求 pH}

\begin{enumerate}[itemsep=2pt]
\item 确定溶液里的 \chem{H^+} 浓度（强酸稀溶液里可近似认为等于酸的浓度）。
\item 取常用对数：$c = 10^{-x}$ 时，$\lg c = -x$。
\item 加负号：$\mathrm{pH} = -(-x) = x$。
\end{enumerate}

例题：某柠檬汁 $c(\chem{H^+}) = 10^{-2.3}\ \ud{}{mol/L}$，pH 就是 2.3。注意换算方向：pH 每\textbf{减小} 1，\chem{H^+} 浓度\textbf{增大} 10 倍。柠檬汁（pH 约 2）的氢离子浓度大约是纯净水（pH 约 7）的 $10^5$ 倍——十万倍，不是“差 5 个点”那么温柔。

\begin{table}[htbp]
\centering
\caption{身边常见液体的 pH（约值，$25\,^{\circ}\mathrm{C}$）}
\begin{tabular}{llll}
\toprule
液体 & pH 约值 & 液体 & pH 约值 \\
\midrule
胃液 & 1.5--2 & 牛奶 & 6.5 \\
柠檬汁 & 2--2.5 & 纯净水 & 7 \\
食醋 & 2.5--3 & 血液 & 7.35--7.45 \\
橙汁 & 3.5 & 海水 & 8.1 \\
酸雨（临界） & 5.6 以下 & 小苏打溶液 & 8--9 \\
雨水（正常） & 5.6 & 肥皂水 & 9--10 \\
\bottomrule
\end{tabular}
\end{table}

这张表有两处值得停一停。一是正常雨水本来就是弱酸性的（空气中的 \chem{CO_2} 溶于水生成碳酸），所以“酸雨”的定义不是“pH 小于 7”，而是“pH 小于 5.6”——判断一件事要先知道天然基线在哪里。二是血液的 pH 被身体死死守在 7.35 到 7.45 之间，靠碳酸氢根等缓冲系统日夜纠偏（第 60 章讲的稳态）。“喝碱性水改变体质”这类说法，等于宣称一杯水能打赢全身这套精密的缓冲系统——剂量和证据都对不上。

\begin{wuqu}
“pH = 7 就是中性，永远如此。”——只在 $25\,^{\circ}\mathrm{C}$ 的纯水里成立。中性的定义是 $c(\chem{H^+}) = c(\chem{OH^-})$，而水的电离程度随温度变化：$100\,^{\circ}\mathrm{C}$ 时纯水 pH 约 6.1，但它依然是中性的，因为两种离子一样多。pH 刻度是温度相关的标尺，7 只是常温下的中点，不是宇宙的常数。把“pH 小于 7”直接等同于“有害”更是错上加错——你的胃液 pH 约 2，是消化工作的正常状态。
\end{wuqu}

反着算也一样常用：知道 pH 求浓度，$c(\chem{H^+}) = 10^{-\mathrm{pH}}$。pH 4.5 的酸雨，\chem{H^+} 浓度约 $10^{-4.5}\ \ud{}{mol/L}$，即 $3\times10^{-5}\ \ud{}{mol/L}$。环境报告里的 pH 数字，就这样还原成真实的粒子浓度。

\section{把工具放回箱子}

这套符号系统的全部秘密，其实就一句话：\textbf{符号是账本的速写，而账本背后永远是真实粒子的排队和能量的流动}。每当你读一个式子读到发懵，就回到那五个贯穿问题：它由什么组成？怎样连接？什么能量规则决定它能不能变？变多快？在生命里怎样被使用？符号答不出的问题，正文每一章都在答。

其他几只“急救箱”也在这里备着：周期表的趋势查附录 B，分子模型的画法与读法查附录 C，遗传信息的术语查附录 D。遇到算不对的题目，附录 F 有提示和解答。

\section*{练一练}

\begin{sikaoti}
\item 数原子：\chem{2(NH_4)_2SO_4}（化肥硫酸铵）里，氮、氢、硫、氧各有几个原子？画出你的“认领下标”过程。
\item 配平并检验：\chem{Fe} + \chem{O_2} $\rightarrow$ \chem{Fe_3O_4}（铁丝在氧气中燃烧）。配平后把两边的铁、氧原子各数一遍，确认账本平了。
\item \chem{Na_2CO_3} 的摩尔质量是多少？\ud{53}{g} 碳酸钠是多少摩尔、约含多少个 \chem{Na^+}？（小心：每 1 个 \chem{Na_2CO_3} 给出 2 个 \chem{Na^+}。）
\item 把 \ud{4.0}{g} 氢氧化钠（\chem{NaOH}，摩尔质量 \ud{40}{g/mol}）配成 \ud{250}{mL} 溶液，物质的量浓度是多少？取其中 \ud{50}{mL} 加水稀释到 \ud{200}{mL}，新浓度又是多少？
\item 两种溶液，甲的 pH 为 3，乙的 pH 为 5。哪种 \chem{H^+} 浓度高？高多少倍？如果乙就是 $60\,^{\circ}\mathrm{C}$ 的纯水，能据此说甲“酸性强、乙是碱”吗？说说理由。
\item 任选一个本章的生物分子记号（葡萄糖、氨基酸通式、ATP、DNA 碱基），用自己的话写一段“这个记号压缩了什么信息、又隐藏了什么信息”的说明，当作你给同学讲解的讲稿。
\end{sikaoti}
